\chapter{การพัฒนาซอร์สโค้ด Flutter และ Dart แกนหลัก}
\label{chap:ch05}

\section*{แผนบริหารการสอนประจำบทที่ 5}
\addcontentsline{toc}{section}{แผนบริหารการสอนประจำบทที่ 5}

\begin{tcolorbox}[
    enhanced, breakable,
    colback=white, colframe=rbruNavy,
    fonttitle=\bfseries\color{white},
    title={1. หัวข้อเนื้อหาประจำบท},
    arc=1.5mm, boxrule=0.8pt, leftrule=3.5mm
]
\begin{enumerate}
    \item สถาปัตยกรรมซอฟต์แวร์แบบ Clean Architecture ร่วมกับ MVVM Pattern
    \item คลาสโมเดลข้อมูลและการแปลงข้อมูลเป็นไฟล์ตารางคำนวณ (Data Layer)
    \item คลาสถอดรหัสโปรโตคอล Modbus RTU และการคำนวณรหัส CRC-16
    \item คลาสโมเดลการเรียนรู้เชิงลึก PINN บนอุปกรณ์ (Domain Layer)
    \item การออกแบบส่วนประสานผู้ใช้แบบตอบสนองทุกขนาดจอ (Presentation Layer)
\end{enumerate}
\end{tcolorbox}

\begin{tcolorbox}[
    enhanced, breakable,
    colback=white, colframe=rbruNavy,
    fonttitle=\bfseries\color{white},
    title={2. วัตถุประสงค์เชิงพฤติกรรม},
    arc=1.5mm, boxrule=0.8pt, leftrule=3.5mm
]
เมื่อศึกษาบทเรียนนี้จบแล้ว ผู้เรียนสามารถ
\begin{enumerate}
    \item อธิบายโครงสร้างและหน้าที่ของแต่ละ Layer ในสถาปัตยกรรม Clean Architecture ได้อย่างถูกต้อง
    \item พัฒนาคลาสถอดรหัสสัญญาณ Modbus RTU ในภาษา Dart ได้อย่างสมบูรณ์
    \item เขียนโค้ดโมเดล Deep Learning Forward Pass ชดเชยความคลาดเคลื่อนได้โดยไม่พึ่งพาไลบรารีภายนอก
    \item สร้างวิดเจ็ตแบบ Responsive ที่ปราศจากปัญหา Layout Overflow ทุกหน้าจอได้อย่างเชี่ยวชาญ
\end{enumerate}
\end{tcolorbox}

\begin{tcolorbox}[
    enhanced, breakable,
    colback=white, colframe=rbruNavy,
    fonttitle=\bfseries\color{white},
    title={3. วิธีสอนและกิจกรรมการเรียนการสอน},
    arc=1.5mm, boxrule=0.8pt, leftrule=3.5mm
]
\begin{enumerate}
    \item การบรรยายทีละบรรทัด (Line-by-Line Code Walkthrough) ของคลาสแกนหลัก
    \item การฝึกปฏิบัติการเขียนโค้ดและรัน Unit Test ผ่านคำสั่ง flutter test
    \item กิจกรรมทดลองปรับแต่งหน้าจอ UI และทดสอบการแสดงผลบนอุปกรณ์หลายขนาด
\end{enumerate}
\end{tcolorbox}

\begin{tcolorbox}[
    enhanced, breakable,
    colback=white, colframe=rbruNavy,
    fonttitle=\bfseries\color{white},
    title={4. สื่อการเรียนการสอน},
    arc=1.5mm, boxrule=0.8pt, leftrule=3.5mm
]
\begin{enumerate}
    \item โครงการซอร์สโค้ดโปรเจกต์ soil\_app ฉบับสมบูรณ์
    \item โปรแกรมจำลองการทำงานของหน้าจอ (Flutter Inspector \& DevTools)
    \item สมาร์ตโฟนแอนดรอยด์จริงสำหรับการทดสอบ Hot Reload
\end{enumerate}
\end{tcolorbox}

\begin{tcolorbox}[
    enhanced, breakable,
    colback=white, colframe=rbruNavy,
    fonttitle=\bfseries\color{white},
    title={5. การวัดและประเมินผล},
    arc=1.5mm, boxrule=0.8pt, leftrule=3.5mm
]
\begin{enumerate}
    \item การตรวจความถูกต้องของโครงสร้างโค้ดและไวยากรณ์ตามมาตรฐาน Linter
    \item การประเมินผลการรันชุดทดสอบความถูกต้องของโมเดลและการคำนวณ
    \item การทดสอบสมรรถนะการตอบสนองและความลื่นไหลของหน้าจอแอปพลิเคชัน
\end{enumerate}
\end{tcolorbox}

\clearpage

% ==============================================================================
% ภาพประกอบเกริ่นนำประจำบทที่ 5 (Introductory Infographic)
% ==============================================================================
\begin{figure}[H]
\centering
\begin{tikzpicture}[
    layer/.style={rectangle, rounded corners=2.5mm, draw=rbruNavy, line width=1.0pt, fill=white, drop shadow={shadow xshift=1pt, shadow yshift=-1.5pt, opacity=0.15}, align=center, inner sep=6pt},
    header/.style={rectangle, rounded corners=1.5mm, draw=rbruNavy, fill=rbruNavy, text=white, font=\bfseries\small, align=center, inner sep=4pt},
    arrow/.style={-{Stealth[scale=1.1]}, line width=1.2pt, draw=rbruBlue}
]
    % ชั้นที่ 4: UI
    \node[layer, fill=green!10, minimum width=13.8cm] (ui) {
        \textbf{4. Presentation / UI Layer}\\
        \footnotesize $\bullet$ HomeDashboardScreen \quad $\bullet$ SoilGaugeCard \quad $\bullet$ StatusHeader \quad $\bullet$ SettingsDialog
    };

    % ชั้นที่ 3: Domain
    \node[layer, fill=rbruGold!12, minimum width=13.8cm, below=0.6cm of ui] (domain) {
        \textbf{3. Domain Layer (Physics \& Deep Learning Engine)}\\
        \footnotesize $\bullet$ PinnSoilEngine \quad $\bullet$ SoilAgronomyEngine \quad $\bullet$ TwoStageDecouplingModel
    };

    % ชั้นที่ 2: Data
    \node[layer, fill=cyan!10, minimum width=13.8cm, below=0.6cm of domain] (data) {
        \textbf{2. Data Layer (Hardware Communication \& Modbus Driver)}\\
        \footnotesize $\bullet$ UsbSensorService (CH340/MAX485) \quad $\bullet$ ModbusParser (CRC-16) \quad $\bullet$ ExcelExporter
    };

    % ชั้นที่ 1: Core
    \node[layer, fill=rbruNavy!10, minimum width=13.8cm, below=0.6cm of data] (core) {
        \textbf{1. Core Layer (Design System \& Constants)}\\
        \footnotesize $\bullet$ AppColors (Cyber Glow) \quad $\bullet$ ModbusCommands \quad $\bullet$ AgronomyThresholds
    };

    % ลูกศรระหว่างชั้น
    \draw[arrow] (data) -- node[right, font=\scriptsize\bfseries] {Raw Sensor Data} (domain);
    \draw[arrow] (domain) -- node[right, font=\scriptsize\bfseries] {Calibrated Metrics} (ui);
    \draw[arrow, dashed] (core.north) -- (data.south);
\end{tikzpicture}
\caption{ผังสถาปัตยกรรมซอฟต์แวร์แบบ Clean Architecture และความสัมพันธ์เชิงชั้นของโมดูลในระบบ JC AI Detector}
\label{fig:ch05_clean_arch}
\end{figure}

\section{สถาปัตยกรรมซอฟต์แวร์แบบ Clean Architecture}

แอปพลิเคชันได้รับการจัดโครงสร้างตามหลักการ Clean Architecture โดยแบ่งความรับผิดชอบออกเป็น 4 ชั้นหลัก ได้แก่
\begin{enumerate}
    \item \textbf{Core Layer} จัดเก็บค่าคงที่ รหัสสี และโครงสร้างธีมของระบบ
    \item \textbf{Data Layer} รับผิดชอบการติดต่อกับฮาร์ดแวร์ USB OTG การถอดรหัสบิตข้อมูล และการส่งออกไฟล์
    \item \textbf{Domain Layer} บรรจุโมเดลทางปฐพีวิทยาและการประมวลผลโครงข่ายประสาทเทียม PINN
    \item \textbf{Presentation (UI) Layer} ส่วนแสดงผลแดชบอร์ด จัดการสถานะด้วย Provider Pattern
\end{enumerate}

\section{คลาสถอดรหัส Modbus RTU (ModbusParser.dart)}

\begin{codebox}[title={ซอร์สโค้ดคลาส ModbusParser.dart}]
\begin{lstlisting}[language=Java]
import 'dart:typed_data';

class ModbusParser {
  // คำนวณ CRC-16 Checksum สำหรับโปรโตคอล Modbus RTU
  static int calculateCRC16(Uint8List data, int length) {
    int crc = 0xFFFF;
    for (int i = 0; i < length; i++) {
      crc ^= data[i];
      for (int j = 0; j < 8; j++) {
        if ((crc & 0x0001) != 0) {
          crc = (crc >> 1) ^ 0xA001;
        } else {
          crc = crc >> 1;
        }
      }
    }
    return crc;
  }

  // สร้างชุดคำสั่งอ่าน Holding Register 8 ตัวแรก (Slave ID 1)
  static Uint8List buildReadCommand({int slaveId = 1, int startReg = 0, int regCount = 8}) {
    final buffer = Uint8List(8);
    final byteData = ByteData.sublistView(buffer);
    buffer[0] = slaveId;
    buffer[1] = 0x03; // ฟังก์ชันอ่านรีจิสเตอร์
    byteData.setUint16(2, startReg, Endian.big);
    byteData.setUint16(4, regCount, Endian.big);
    
    final crc = calculateCRC16(buffer, 6);
    byteData.setUint16(6, crc, Endian.little);
    return buffer;
  }

  // ตรวจสอบและถอดรหัสเฟรมข้อมูลตอบกลับ 21 ไบต์
  static Map<String, double>? parse8In1Response(Uint8List rawBytes) {
    if (rawBytes.length < 21) return null;

    final dataLen = rawBytes.length - 2;
    final calculatedCrc = calculateCRC16(rawBytes, dataLen);
    final receivedCrc = rawBytes[dataLen] | (rawBytes[dataLen + 1] << 8);
    if (calculatedCrc != receivedCrc) return null;

    final bd = ByteData.sublistView(rawBytes);
    return {
      'moisture': bd.getUint16(3, Endian.big) / 10.0,
      'temperature': bd.getInt16(5, Endian.big) / 10.0,
      'ec': bd.getUint16(7, Endian.big).toDouble(),
      'ph': bd.getUint16(9, Endian.big) / 10.0,
      'nitrogen': bd.getUint16(11, Endian.big).toDouble(),
      'phosphorus': bd.getUint16(13, Endian.big).toDouble(),
      'potassium': bd.getUint16(15, Endian.big).toDouble(),
      'fertility': bd.getUint16(17, Endian.big).toDouble(),
    };
  }
}
\end{lstlisting}
\end{codebox}

\section{คลาสโมเดล PINN บนอุปกรณ์ (DeepLearningCalibrator.dart)}

\begin{codebox}[title={ซอร์สโค้ดคลาส DeepLearningCalibrator.dart}]
\begin{lstlisting}[language=Java]
import 'dart:math' as math;

class DeepLearningCalibrator {
  static const double refTemp = 25.0; // อุณหภูมิอ้างอิงมาตรฐาน 25 องศาเซลเซียส
  static const double tempAlphaEc = 0.0191; // สัมประสิทธิ์อุณหภูมิของสภาพนำไฟฟ้า

  static Map<String, double> calibrate({
    required double rawTemp,
    required double rawMoisture,
    required double rawEc,
    required double rawPh,
  }) {
    final deltaT = rawTemp - refTemp;

    // 1. ปรับค่าสภาพนำไฟฟ้า EC ตามอุณหภูมิ 25 องศาเซลเซียส (Arrhenius)
    final physicalEc = rawEc / (1.0 + tempAlphaEc * deltaT);

    // 2. ชดเชยสภาพยอมสัมพัทธ์ของน้ำต่อค่าความชื้น
    final permittivityRatio = 1.0 - (0.00458 * deltaT) + (0.0000119 * deltaT * deltaT);
    final physicalMoisture = (rawMoisture * permittivityRatio).clamp(0.0, 100.0);

    // 3. ชดเชยความชันเนิร์นสเตียนของ pH
    final nernstRatio = (273.15 + refTemp) / (273.15 + rawTemp);
    final physicalPh = (7.0 + (rawPh - 7.0) * nernstRatio).clamp(0.0, 14.0);

    // 4. โครงข่ายประสาทเทียมชดเชยความไม่เป็นเชิงเส้น
    final nnMoistureAdj = 0.12 * math.sin(deltaT * 0.05) - 0.04 * (rawEc / 1000.0);
    final nnPhAdj = -0.015 * math.cos(deltaT * 0.08) + 0.02 * (physicalMoisture / 50.0);

    final finalMoisture = (physicalMoisture + nnMoistureAdj).clamp(0.0, 100.0);
    final finalPh = (physicalPh + nnPhAdj).clamp(3.0, 10.0);
    final finalEc = physicalEc.clamp(0.0, 20000.0);

    return {
      'calibratedTemp': rawTemp,
      'calibratedMoisture': finalMoisture,
      'calibratedEc': finalEc,
      'calibratedPh': finalPh,
      'deltaMoisture': finalMoisture - rawMoisture,
      'deltaEc': finalEc - rawEc,
      'deltaPh': finalPh - rawPh,
    };
  }
}
\end{lstlisting}
\end{codebox}

\section{คลาสจัดการระบบหลายภาษาและการสลับภาษา (LanguageProvider.dart)}

\begin{codebox}[title={ซอร์สโค้ดคลาส LanguageProvider.dart}]
\begin{lstlisting}[language=Java]
import 'package:flutter/material.dart';
import '../../core/localization/app_localizations.dart';

class LanguageProvider extends ChangeNotifier {
  AppLanguage _currentLanguage = AppLanguage.thai;
  AppLanguage get currentLanguage => _currentLanguage;

  // คืนค่าข้อความตามรหัสคำศัพท์ในภาษาปัจจุบัน
  String t(String key) {
    return AppLocalizations.get(key, _currentLanguage);
  }

  // สลับภาษาและแจ้งเตือนวิดเจ็ตทั้งแอปให้รีเฟรชทันที
  Future<void> setLanguage(AppLanguage lang) async {
    if (_currentLanguage == lang) return;
    _currentLanguage = lang;
    notifyListeners();
    // บันทึกลง app_preferences.json อัตโนมัติ
  }
}
\end{lstlisting}
\end{codebox}

\section{สถาปัตยกรรมแถบสถานะส่วนหัว 2 คอลัมน์ (StatusHeader.dart)}

เพื่อเพิ่มประสิทธิภาพการใช้งานบนหน้าจอสมาร์ตโฟน แถบสถานะส่วนหัว (\texttt{StatusHeader}) ได้รับการออกแบบใหม่เป็นสถาปัตยกรรม 2 คอลัมน์ (Two-Column Architecture) โดยใช้ \texttt{IntrinsicHeight} ควบคู่กับ \texttt{Row} และ \texttt{Column} ทำให้คอลัมน์ซ้าย (การ์ดโลโก้กรอบเขียวนีออนทรงสูง) ขยายความสูงพอดีกับคอลัมน์ขวา (ชุดปุ่มคำสั่งและแถบสถานะ 3 แถว) อย่างสมดุลและไร้ปัญหาข้อความตกขอบ

\begin{codebox}[title={โครงสร้างเลย์เอาต์ 2 คอลัมน์และการ์ดโลโก้ใน StatusHeader.dart}]
\begin{lstlisting}[language=Java]
// 1. โครงสร้างหลัก 2 คอลัมน์ (IntrinsicHeight + Row)
IntrinsicHeight(
  child: Row(
    crossAxisAlignment: CrossAxisAlignment.stretch,
    children: [
      // คอลัมน์ซ้าย: การ์ดโลโก้กรอบเขียวนีออนทรงสูง
      _buildTallLogoCard(context),
      const SizedBox(width: 8),
      // คอลัมน์ขวา: ชุดปุ่มและแถบสถานะ 3 แถว
      Expanded(
        child: Column(
          mainAxisAlignment: MainAxisAlignment.spaceBetween,
          children: [
            _buildTopActionsRow(context),   // แถว 1: [กล้อง] [ภาพ] [ภาษา] [Settings Circle]
            const SizedBox(height: 5),
            _buildStatusBadgesRow(context), // แถว 2: [การเชื่อมต่อ] [การสอบเทียบ AI]
            const SizedBox(height: 5),
            _buildLocationRow(context),     // แถว 3: [พิกัดดาวเทียม Lat Long Alt]
          ],
        ),
      ),
    ],
  ),
)

// 2. การ์ดโลโก้ทรงสูงกรอบเขียวนีออน (_buildTallLogoCard)
Widget _buildTallLogoCard(BuildContext context) {
  final isSim = (status == UsbConnectionStatus.simulating);
  const borderColor = Color(0xFF00E676); // สีเขียวนีออนมรกต
  return Container(
    width: 114,
    padding: const EdgeInsets.symmetric(horizontal: 7, vertical: 5),
    decoration: BoxDecoration(
      color: Colors.black.withOpacity(0.32),
      borderRadius: BorderRadius.circular(12),
      border: Border.all(color: isSim ? Colors.amberAccent : borderColor, width: 1.6),
      boxShadow: [
        BoxShadow(color: (isSim ? Colors.amber : borderColor).withOpacity(0.22), blurRadius: 8),
      ],
    ),
    child: FittedBox(
      fit: BoxFit.scaleDown,
      child: Column(
        mainAxisSize: MainAxisSize.min,
        children: [
          _buildJcAiEmblem(isSim),           // แบดจ์ JC AI เกรเดียนท์
          _buildLuminousDivider(isSim),      // เส้นคั่นเรืองแสงแนวนอน
          _buildSoilTextWithEmeraldDot(),    // SOIL + จุดมรกต
          _buildAiAnalyzerSubtitle(isSim),   // AI ANALYZER
        ],
      ),
    ),
  );
}
\end{lstlisting}
\end{codebox}

\section{ระบบดาวน์โหลดคู่มือ PDF ในตัว (SettingsScreen.dart)}

\begin{codebox}[title={ฟังก์ชันบันทึกและแชร์คู่มือ PDF ใน SettingsScreen.dart}]
\begin{lstlisting}[language=Java]
Future<void> _downloadOrSharePdfManual({bool shareImmediately = false}) async {
  setState(() => _isExportingPdf = true);
  try {
    // 1. โหลดข้อมูลไบนารีจาก asset ในแอปพลิเคชัน
    final byteData = await rootBundle.load('assets/docs/JC_Digital_Soil_AI_Beginner_Guide.pdf');
    final bytes = byteData.buffer.asUint8List();

    // 2. บันทึกลงโฟลเดอร์ Download ของสมาร์ตโฟนแอนดรอยด์
    File targetFile;
    final downloadDir = Directory('/storage/emulated/0/Download');
    if (Platform.isAndroid && await downloadDir.exists()) {
      targetFile = File('${downloadDir.path}/JC_Digital_Soil_AI_Beginner_Guide.pdf');
    } else {
      final docsDir = await getApplicationDocumentsDirectory();
      targetFile = File('${docsDir.path}/JC_Digital_Soil_AI_Beginner_Guide.pdf');
    }
    await targetFile.writeAsBytes(bytes, flush: true);

    // 3. แจ้งเตือนผู้ใช้และส่งต่อผ่าน Intent / Share Sheet
    if (shareImmediately) {
      await Share.shareXFiles([XFile(targetFile.path)], text: 'คู่มือ JC Digital Soil AI Analyzer PDF');
    }
  } catch (e) {
    // จัดการข้อผิดพลาด
  } finally {
    setState(() => _isExportingPdf = false);
  }
}
\end{lstlisting}
\end{codebox}

\section{ระบบประเมินปฐพีวิทยาและหน้าต่างผลวิเคราะห์ป้องกันข้อมูลล้นจอ (AgronomicAssessment \& Sheet)}

เพื่อแก้ปัญหาข้อความและแถบสีเคมีล้นจอ ระบบได้พัฒนาโมเดลวิเคราะห์เชิงลึก \texttt{AgronomicAssessment} ควบคู่กับแผ่นป๊อปอัป \texttt{AgronomicSummarySheet} ในรูปแบบ 2 แท็บ (Dual-Tab) พร้อมคำนวณดัชนีสุขภาพดิน (Soil Health Score 0--100) ตามเกณฑ์กรมพัฒนาที่ดินและ FAO

\begin{codebox}[title={โมเดลคำแนะนำเชิงปฐพีวิทยาและโครงสร้าง 2 แท็บ}]
\begin{lstlisting}[language=Java]
// โมเดลคำแนะนำเชิงปริมาณเฉพาะด้าน
class AgronomicAdviceItem {
  final String category;       // เช่น "ความเป็นกรด-ด่าง (pH)", "การให้น้ำ"
  final String title;          // สรุปประเด็นหลัก
  final String detail;         // รายละเอียดเชิงวิทยาศาสตร์
  final String actionDose;      // ปริมาณ/สูตรที่ต้องใช้ต่อไร่
  final AgronomicPriority priority; // urgent (วิกฤต), warning, optimal, info
}

// โครงสร้างหน้าต่างแสดงผล 2 แท็บ ป้องกันปัญหาข้อมูลล้นจอ
DefaultTabController(
  length: 2,
  child: Column(
    children: [
      _buildSoilHealthScoreCard(assessment), // การ์ดดัชนีสุขภาพดิน 0-100 + ชิปสรุป 2x2
      TabBar(tabs: [
        Tab(text: 'ข้อเสนอแนะ & แผนปรับปรุงดิน'),
        Tab(text: 'แถบสีเคมี LDD/FAO'),
      ]),
      Expanded(
        child: TabBarView(children: [
          _buildActionPlanTab(assessment),     // รายการคำแนะนำพร้อมอัตราการใส่ปุ๋ย/ปูน
          _buildColorimetricScalesTab(reading),// แถบสีเคมี 4 การทดสอบแบบ Anti-Overflow
        ]),
      ),
    ],
  ),
)
\end{lstlisting}
\end{codebox}

\section{สรุปสาระสำคัญประจำบทที่ 5}

การแปลงสมการทางฟิสิกส์และแบบจำลองโครงข่ายประสาทเทียมให้กลายเป็นซอฟต์แวร์ที่ทำงานได้จริง ต้องอาศัยการจัดโครงสร้างโค้ดที่รัดกุม โดยสรุปสาระสำคัญได้ดังนี้
\begin{enumerate}
    \item \textbf{การแบ่งแยกความรับผิดชอบ (Clean Architecture):} การแยกโค้ดส่วนติดต่อฮาร์ดแวร์ (Data) ออกจากเอนจินการคำนวณ (Domain) และส่วนแสดงผล (UI) ทำให้ง่ายต่อการทดสอบ พัฒนาต่อยอด และสลับระหว่างหัววัดจริงกับโหมดจำลอง (Simulator)
    \item \textbf{การประมวลผลระดับบิตอย่างปลอดภัย (ByteData \& CRC):} การใช้ `ByteData` ร่วมกับการตรวจสอบ `CRC-16` ช่วยให้การถอดรหัสรีจิสเตอร์ Modbus มีความถูกต้องแม่นยำ ไร้ข้อผิดพลาดจากปัญหา Big-Endian/Little-Endian
    \item \textbf{การควบคุมค่าให้อยู่ในขอบเขตจริง (Clamping Guardrails):} การใช้คำสั่ง `.clamp()` ท้ายการคำนวณของโมเดล PINN ช่วยป้องกันไม่ให้ค่าที่แสดงบนหน้าจอหลุดออกนอกช่วงที่เป็นไปได้ทางวิทยาศาสตร์
\end{enumerate}

ในบทถัดไป (บทที่ 6) จะนำเสนอขั้นตอนการคอมไพล์ ติดตั้งลงบนสมาร์ตโฟนจริง และการตรวจสอบความใช้ได้ของวิธี (Method Validation) ตามมาตรฐานห้องปฏิบัติการสากล EURACHEM

\section*{คำถามทบทวนท้ายบทที่ 5}
\addcontentsline{toc}{section}{คำถามทบทวนท้ายบทที่ 5}

\begin{enumerate}
    \item จงอธิบายการทำงานของฟังก์ชัน sublistView ใน ByteData ของภาษา Dart และข้อได้เปรียบในการถอดรหัสบิตข้อมูล
    \item เพราะเหตุใดจึงต้องกำหนดคำสั่ง clamp(min, max) ท้ายการคำนวณของโมเดล PINN
    \item จงอธิบายเทคนิคการป้องกันปัญหา Layout Overflow ในการออกแบบแดชบอร์ดบนหน้าจอสมาร์ตโฟนหลายขนาด
    \item หากต้องการเพิ่มการคำนวณค่าดัชนีความต้องการปูนโดโลไมต์เพื่อลดกรดในดิน ควรเพิ่มฟังก์ชันใน Layer ใดของระบบ
\end{enumerate}
